\documentclass{article}
\usepackage{mathtools} % loads amsmath
\usepackage{amssymb}
\usepackage{amsthm}
\usepackage{hyperref}
\usepackage{cleveref} % must come after hyperref
\usepackage{thmtools} % \declaretheorem: lets \cref name shared-counter environments correctly

\declaretheorem[name=Theorem]{theorem}
\declaretheorem[name=Lemma, sibling=theorem]{lemma}
\declaretheorem[name=Corollary, sibling=theorem]{corollary}
\declaretheorem[name=Definition, style=definition, sibling=theorem]{definition}

\title{Project Euler Problem 2: Even Fibonacci Numbers}
\author{Brinhasavlin}
\date{\today}

\begin{document}
\maketitle

\section*{Problem}

Each new term in the Fibonacci sequence is generated by adding the previous two terms. By starting with $1$ and $2$, the
first $10$ terms will be:
\[1, 2, 3, 5, 8, 13, 21, 34, 55, 89, \dots\]

By considering the terms in the Fibonacci sequence whose values do not exceed four million, find the sum of the
even-valued terms.

\section*{Proof}

\begin{definition}\label{def:fib}
  The $n$th Fibonacci number $\operatorname{F}(n)$ is given by the following relation:
\begin{align*}
  \operatorname{F}(n) &= \operatorname{F}(n-1) + \operatorname{F}(n-2)  && \text{for $n > 2$} \\
  \operatorname{F}(1) &= 1 \\
  \operatorname{F}(2) &= 2
\end{align*}
\end{definition}

One interesting property that arises from this definition is that starting from $\operatorname{F}(2)$, every third
Fibonacci number is even.
\begin{lemma}\label{lem:parity}
  For $n \ge 0$, $\operatorname{F}(3n + 1)$ is odd, $\operatorname{F}(3n + 2)$ is even, and $\operatorname{F}(3n + 3)$
  is odd.
\end{lemma}
\begin{proof}
  We proceed by induction on $n$.

  \textbf{Base case} ($n = 0$).
  \Cref{def:fib} gives this almost immediately: $\operatorname{F}(1) = 1$ is odd,
  $\operatorname{F}(2) = 2$ is even, and
  \[\operatorname{F}(3) = \operatorname{F}(1) + \operatorname{F}(2) = 3\] is odd.
  Thus, all three statements hold for $n=0$.

  \textbf{Inductive step}.
  Suppose that for some $n \ge 0$, all three statements hold. We show they also hold for $n+1$. First,
  \[\operatorname{F}(3n + 4) = \operatorname{F}(3n+3) + \operatorname{F}(3n+2)\]
  We know that $\operatorname{F}(3n + 3)$ is odd and $\operatorname{F}(3n+2)$ is even. Thus, $\operatorname{F}(3n + 4)$
  must be odd, since the sum of an odd and an even number is odd.

  Next,
  \[\operatorname{F}(3n + 5) = \operatorname{F}(3n+4) + \operatorname{F}(3n+3)\]
  We just showed that $\operatorname{F}(3n + 4)$ is odd, and by hypothesis $\operatorname{F}(3n+3)$ is odd. Thus,
  $\operatorname{F}(3n + 5)$ must be even, since the sum of two odd numbers is even.

  Finally,
  \[\operatorname{F}(3n + 6) = \operatorname{F}(3n+5) + \operatorname{F}(3n+4)\]
  We have shown that $\operatorname{F}(3n + 5)$ is even and $\operatorname{F}(3n+4)$ is odd. Thus,
  $\operatorname{F}(3n + 6)$ must be odd, since the sum of an even and an odd number is odd.

  Since $3(n+1) + 1 = 3n + 4$, $3(n+1) + 2 = 3n + 5$, and $3(n+1) + 3 = 3n + 6$, all three statements hold for $n + 1$.
  By induction, they hold for all $n \ge 0$.
\end{proof}

\begin{corollary}\label{cor:even}
  For $n \ge 1$, $\operatorname{F}(n)$ is even if and only if $n \equiv 2 \pmod{3}$.
\end{corollary}
\begin{proof}
  Every $n \ge 1$ can be written in exactly one of the forms $3k + 1$, $3k + 2$, or $3k + 3$ for some $k \ge 0$. By
  \cref{lem:parity}, $\operatorname{F}(n)$ is even in the second case and odd in the other two. Thus,
  $\operatorname{F}(n)$ is even exactly when $n = 3k + 2$, that is, when $n \equiv 2 \pmod{3}$.
\end{proof}


Following \Cref{cor:even}, we can develop a recurrence relation that generates the even Fibonacci numbers directly,
skipping the odd ones.

\begin{lemma}\label{lem:relation}
  Let $\operatorname{E}(n)$ be the $n$th even Fibonacci number.
  \begin{align*}
    \operatorname{E}(n) &= 4\operatorname{E}(n-1) + \operatorname{E}(n-2)  &&\text{for $n > 2$} \\
    \operatorname{E}(1) &= 2 \\
    \operatorname{E}(2) &= 8
  \end{align*}
\end{lemma}
\begin{proof}
  By \Cref{cor:even}, the even terms are exactly $\operatorname{F}(2), \operatorname{F}(5), \operatorname{F}(8), \dots$,
  so $\operatorname{E}(n) = \operatorname{F}(3n-1)$. This gives the base cases directly:
  $\operatorname{E}(1) = \operatorname{F}(2) = 2$ and $\operatorname{E}(2) = \operatorname{F}(5) = 8$.

  Now let $n > 2$ and $m = 3n - 4$. Then $\operatorname{E}(n) = \operatorname{F}(m+3)$,
  $\operatorname{E}(n-1) = \operatorname{F}(m)$, and $\operatorname{E}(n-2) = \operatorname{F}(m-3)$.

  With this formulation, it suffices to show that
  $\operatorname{F}(m+3) = 4\operatorname{F}(m) + \operatorname{F}(m-3)$.

  \begin{align*}
    \operatorname{F}(m+3) &= \operatorname{F}(m + 2) + \operatorname{F}(m+1)  && \text{by definition} \\
                          &=  2\operatorname{F}(m + 1) + \operatorname{F}(m) && \text{expanding $\operatorname{F}(m+2)$} \\
                          &= 3\operatorname{F}(m) + 2\operatorname{F}(m-1) && \text{expanding $\operatorname{F}(m+1)$}\\
                          &= 3\operatorname{F}(m) + \operatorname{F}(m-1) + \operatorname{F}(m-2)
                                  && \text{expanding one copy of $\operatorname{F}(m-1)$} \\
                          &\qquad + \operatorname{F}(m-3) \\
                          &= 4\operatorname{F}(m) + \operatorname{F}(m-3)
                                  && \text{regrouping by \Cref{def:fib}}
  \end{align*}
  Substituting back gives $\operatorname{E}(n) = 4\operatorname{E}(n-1) + \operatorname{E}(n-2)$, as required.
\end{proof}

\begin{theorem}\label{thm:answer}
  The sum of the even Fibonacci numbers not exceeding $4{,}000{,}000$ is $4{,}613{,}732$.
\end{theorem}
\begin{proof}
  By \Cref{cor:even}, the even Fibonacci numbers are exactly the terms $\operatorname{E}(k)$. The sequence
  $\operatorname{E}$ is strictly increasing: $\operatorname{E}(2) > \operatorname{E}(1) > 0$, and for $k > 2$,
  $\operatorname{E}(k) = 4\operatorname{E}(k-1) + \operatorname{E}(k-2) > \operatorname{E}(k-1)$. So we can generate the
  terms in order with \Cref{lem:relation} and stop at the first one that exceeds $4{,}000{,}000$.

  \begin{center}
    \begin{tabular}{rrr}
      $k$ & $\operatorname{E}(k)$ & $\sum_{i=1}^{k} \operatorname{E}(i)$ \\
      \hline
      1  & $2$            & $2$           \\
      2  & $8$            & $10$          \\
      3  & $34$           & $44$          \\
      4  & $144$          & $188$         \\
      5  & $610$          & $798$         \\
      6  & $2{,}584$      & $3{,}382$     \\
      7  & $10{,}946$     & $14{,}328$    \\
      8  & $46{,}368$     & $60{,}696$    \\
      9  & $196{,}418$    & $257{,}114$   \\
      10 & $832{,}040$    & $1{,}089{,}154$ \\
      11 & $3{,}524{,}578$ & $4{,}613{,}732$ \\
      \hline
      12 & $14{,}930{,}352$ & ---         \\
    \end{tabular}
  \end{center}

  Since $\operatorname{E}(12) > 4{,}000{,}000$, the even terms not exceeding $4{,}000{,}000$ are
  $\operatorname{E}(1)$ through $\operatorname{E}(11)$, and their sum is $4{,}613{,}732$.
\end{proof}

Generating only the even terms takes 11 steps, whereas the $32$ Fibonacci terms up to $4{,}000{,}000$ would each need to
be generated and checked for parity.

\end{document}
